\newglossaryentry{Linux} % Etiqueta que se usa para hacer referencia a la entrada del glosario
{
	name=linux, % Palabra o frase exacta que aparece en el glosario
    plural=linuces,
	description={is a generic term referring to the family of Unix-like
		computer operating systems that use the Linux kernel}	
}

\newacronym{pq}{PQ}{Page Quality}

%\newglossaryentry{application programming interface} % Etiqueta que se usa para hacer referencia a la entrada del glosario
%{
%    name=application programming interface, % Palabra o frase exacta que aparece en el glosario y en el cuerpo del texto.
%    plural=application programming interfaces,
%    description={es una pieza de código que permite a dos aplicaciones comunicarse entre sí para compartir información y funcionalidades. Se usan generalmente en bibliotecas de programación.[}	
%}

\newacronym[plural=APIs]
{API} % etiqueta para hacer referencia
{API} % entrada del glosario y cómo se muestra en el texto
{Una API (\textit{Application Programming Interface} / Interfaz de Programación de Aplicaciones) es una pieza de código que permite a dos aplicaciones comunicarse entre sí para compartir información y funcionalidades. Se usan generalmente en bibliotecas de programación.}


\newglossaryentry{computer}{name=computer, description={A programmable electronic device}}

% {label}{}
\newglossaryentry{modeling language}{
    name={modeling language}, 
    description={Lenguaje con reglas definidas para describir la estructura de un sistema de forma estandarizada.}
}